\chapter{ОХРАНА ТРУДА}

\section{Требования безопасности при работе с вычислительной техникой}

Перед включением оборудования необходимо визуально проверить рабочее место. Провода не должны быть пережаты ножками стола или кресла, блоки питания не должны находиться в местах возможного перегрева, а розетки должны быть исправны. При работе с ноутбуком нужно использовать оригинальный или совместимый блок питания, рассчитанный на требуемую мощность.


Нельзя выполнять влажную уборку оборудования, когда оно подключено к сети. Жидкости должны находиться на безопасном расстоянии от клавиатуры, системного блока и сетевых фильтров. Попадание жидкости может привести не только к повреждению устройства, но и к короткому замыканию.


При возникновении запаха гари, искрения, сильного шума вентилятора или нестабильной работы компьютера необходимо прекратить работу. Оборудование следует выключить штатным способом, если это безопасно, отключить питание и сообщить ответственному лицу. Самостоятельный ремонт электрической части без квалификации недопустим.


Работа программиста связана с длительным использованием персонального компьютера, монитора, клавиатуры, мыши, сетевого оборудования и периферийных устройств. Перед началом работы необходимо убедиться в исправности оборудования, целостности кабелей, надежности подключения питания и отсутствии внешних повреждений корпуса. Запрещается работать с устройствами, имеющими поврежденные провода, признаки перегрева, запах гари, нестабильное питание или открытые токоведущие элементы.


Электробезопасность обеспечивается использованием исправных розеток, сетевых фильтров и оборудования с заводскими блоками питания. Не допускается подключение большого числа устройств к одной перегруженной линии, самостоятельный ремонт блока питания, разборка монитора или системного блока при включенном питании. При обнаружении неисправности необходимо выключить устройство и обратиться к ответственному специалисту.


Пожарная безопасность требует поддержания рабочего места в порядке. Бумага, ткань и другие легковоспламеняющиеся материалы не должны находиться вплотную к нагревающимся элементам. Вентиляционные отверстия системного блока и ноутбука нельзя закрывать. После завершения работы следует корректно выключить оборудование или перевести его в предусмотренный режим, если это соответствует правилам помещения.


При работе за компьютером необходимо соблюдать режим труда и отдыха. Непрерывная длительная работа без перерывов повышает риск зрительного напряжения, утомления мышц спины, шеи, кистей и снижения концентрации. Поэтому рабочий процесс должен включать короткие паузы, смену позы и упражнения для глаз.


\section{Организация рабочего места программиста}

При работе над дипломным проектом программист часто одновременно использует IDE, терминал, браузер, документацию и мессенджер. Поэтому важно организовать рабочее пространство так, чтобы не приходилось постоянно наклоняться к экрану или искать нужное окно. Если используется один монитор, рекомендуется разделить рабочие области программно. Если используются два монитора, основной экран должен находиться прямо перед пользователем.


Клавиатура должна располагаться так, чтобы кисти не были постоянно согнуты вверх. Мышь должна находиться рядом с клавиатурой, чтобы рука не уходила далеко в сторону. При длительной работе желательно использовать коврик с удобной поверхностью и регулярно менять нагрузку между мышью и горячими клавишами.


Документы, методичка и записи по диплому должны быть расположены так, чтобы их можно было читать без сильного поворота головы. Если бумажные материалы используются постоянно, удобна подставка для документов рядом с монитором. Это снижает нагрузку на шею и ускоряет работу.


Рабочее место программиста должно обеспечивать устойчивое положение тела, удобный обзор монитора и свободный доступ к клавиатуре и мыши. Стол должен иметь достаточную площадь для размещения монитора, клавиатуры, мыши, документов и других необходимых предметов. Поверхность стола должна быть устойчивой и не создавать бликов.


Кресло должно поддерживать спину и позволять регулировать высоту посадки. Спина должна опираться на спинку кресла, плечи должны быть расслаблены, локти располагаться примерно под прямым углом, а стопы — устойчиво стоять на полу или подставке. Неправильная посадка приводит к статической нагрузке и повышает риск болей в спине и шее.


Монитор следует располагать на расстоянии примерно 50-70 см от глаз. Верхняя часть экрана должна находиться примерно на уровне глаз или немного ниже. Яркость и контрастность необходимо настроить так, чтобы изображение было четким, но не вызывало чрезмерного напряжения. При работе с кодом и документацией важно избегать мелкого шрифта и чрезмерно ярких цветовых схем.


Освещение должно быть равномерным. Недопустимы сильные блики на экране и резкий контраст между монитором и окружающей средой. При естественном освещении монитор желательно располагать боком к окну. Искусственный источник света должен освещать рабочую область, не направляя яркий поток прямо в глаза.


\section{Анализ вредных и опасных факторов}

При разработке программного обеспечения вредные факторы часто проявляются постепенно. Пользователь может не сразу заметить, что сидит в неудобной позе, работает при слишком высокой яркости или долго не делает перерыв. Поэтому профилактические меры должны быть регулярными, а не только реакцией на уже возникшее утомление.


Зрительное напряжение усиливается при работе с мелким текстом кода, особенно если одновременно используются темная IDE и яркие веб-страницы. При подготовке дипломной записки добавляется работа с большими текстовыми документами, таблицами и рисунками. Это требует более внимательной настройки масштаба и освещения.


Психоэмоциональная нагрузка возрастает при отладке сложных ошибок. Например, если backend не запускается из-за базы данных или frontend не может подключиться к WebSocket, разработчик может долго находиться в состоянии концентрации и напряжения. Чтобы снизить этот фактор, полезно фиксировать ошибки, проверять гипотезы последовательно и не пытаться исправлять несколько причин одновременно.


Основным вредным фактором при работе программиста является зрительное напряжение. Оно возникает из-за длительной фиксации взгляда на мониторе, мелких деталей интерфейса, высокой яркости, бликов и недостаточных перерывов. Последствиями могут быть сухость глаз, головная боль, снижение концентрации и утомляемость.


Второй фактор — статическая нагрузка. При длительной работе в одной позе мышцы спины, шеи и плечевого пояса находятся в напряжении. Это может приводить к боли, нарушению осанки и снижению работоспособности. Нагрузка на кисти и пальцы возникает при продолжительном наборе текста и использовании мыши.


Психоэмоциональное напряжение связано с необходимостью решать сложные логические задачи, исправлять ошибки, соблюдать сроки и поддерживать концентрацию. При разработке дипломного проекта это проявляется в длительном анализе кода, отладке backend/frontend и подготовке документации.


К дополнительным факторам относятся электромагнитное воздействие от оборудования, шум вентиляторов и внешней среды, неудовлетворительный микроклимат, недостаточная вентиляция, неправильная температура и влажность воздуха. В современных условиях эти факторы обычно не превышают допустимые значения, но требуют контроля.


\noindent\scriptsize
\setlength{\tabcolsep}{3pt}
\renewcommand{\arraystretch}{1.18}
\begin{longtable}{|>{\RaggedRight\arraybackslash}p{0.20\linewidth}|>{\RaggedRight\arraybackslash}p{0.28\linewidth}|>{\RaggedRight\arraybackslash}p{0.28\linewidth}|>{\RaggedRight\arraybackslash}p{0.12\linewidth}|}
\caption{Вредные и опасные факторы при работе над проектом}\label{tab:3-1}\\
\hline
Фактор & Причина возникновения & Возможные последствия & Уровень риска\\ \hline
\endfirsthead
\hline
Фактор & Причина возникновения & Возможные последствия & Уровень риска\\ \hline
\endhead
Зрительное напряжение & Длительная работа с кодом, интерфейсами, таблицами, PDF и браузером & Усталость глаз, сухость, головная боль, снижение концентрации & Средний\\ \hline
Статическая нагрузка & Долгое сидение в одной позе, неудобная высота стола или кресла & Боль в спине, шее, плечах, снижение работоспособности & Средний\\ \hline
Нагрузка на кисти & Длительный набор текста, частое использование мыши, неправильное положение запястий & Утомление кистей, дискомфорт, онемение пальцев & Средний\\ \hline
Психоэмоциональное напряжение & Отладка ошибок, сроки сдачи, необходимость одновременно вести код и документацию & Раздражительность, усталость, потеря внимания, снижение качества решений & Средний\\ \hline
Электробезопасность & Поврежденные кабели, перегруженные розетки, нештатные блоки питания & Риск короткого замыкания, повреждения оборудования, травмы & Низкий при соблюдении правил\\ \hline
Пожарная безопасность & Перегрев оборудования, закрытая вентиляция, жидкости и бумага рядом с техникой & Повреждение устройств, задымление, пожарная опасность & Низкий при соблюдении правил\\ \hline
Микроклимат & Душное помещение, высокая температура, недостаточная вентиляция & Быстрая утомляемость, снижение внимания, дискомфорт & Низкий/средний\\ \hline
\end{longtable}
\normalsize


Из таблицы видно, что наиболее значимыми для разработки ChessView являются факторы, связанные с длительной концентрацией и статической работой за компьютером. При соблюдении режима труда и отдыха эти риски остаются управляемыми, но их нельзя игнорировать, так как подготовка дипломного проекта включает не только программирование, но и продолжительную работу с текстом пояснительной записки.


\section{Мероприятия по снижению рисков}

Режим труда и отдыха должен быть заранее запланирован. Для длительной работы за компьютером рекомендуется делать короткие перерывы каждые 45-60 минут. Во время перерыва желательно встать, пройтись, размять плечи, шею и кисти, а также перевести взгляд на удаленные объекты.


Организация рабочего места должна включать контроль температуры и проветривание. Душное помещение снижает концентрацию и повышает утомляемость. Оптимальный микроклимат помогает поддерживать устойчивую работоспособность, особенно при длительной подготовке документации и проверке проекта.


Для снижения нагрузки на кисти можно чередовать ввод текста, чтение документации и работу с мышью. При большом объеме набора текста важно следить, чтобы запястья не опирались на острый край стола. При появлении боли или онемения необходимо сделать перерыв и изменить положение рук.


Для снижения зрительного напряжения необходимо делать регулярные перерывы, переводить взгляд на удаленные объекты, выполнять простую гимнастику для глаз и настраивать яркость монитора под освещение помещения. Текст на экране должен быть достаточно крупным, а рабочие окна — организованы так, чтобы не приходилось постоянно напрягать зрение.


Для уменьшения статической нагрузки следует контролировать посадку, использовать кресло с поддержкой спины, держать клавиатуру и мышь на удобной высоте, делать короткую разминку и менять положение тела. Желательно избегать работы на ноутбуке без подставки и внешней клавиатуры при длительных сессиях.


Кабели питания и периферии должны быть организованы так, чтобы не создавать риск спотыкания или случайного отключения. Рабочее место необходимо поддерживать в чистоте, не размещать жидкости рядом с оборудованием и не закрывать вентиляционные отверстия.


Для снижения психоэмоционального напряжения рекомендуется планировать работу небольшими этапами, фиксировать прогресс, выполнять проверку по чек-листам и избегать длительной непрерывной работы без отдыха. Такой подход особенно важен при подготовке дипломного проекта, где требуется одновременно поддерживать качество кода, документации и демонстрационных материалов.


\noindent\scriptsize
\setlength{\tabcolsep}{3pt}
\renewcommand{\arraystretch}{1.18}
\begin{longtable}{|>{\RaggedRight\arraybackslash}p{0.24\linewidth}|>{\RaggedRight\arraybackslash}p{0.34\linewidth}|>{\RaggedRight\arraybackslash}p{0.24\linewidth}|}
\caption{Профилактические мероприятия при работе за компьютером}\label{tab:3-2}\\
\hline
Направление & Мероприятие & Ожидаемый эффект\\ \hline
\endfirsthead
\hline
Направление & Мероприятие & Ожидаемый эффект\\ \hline
\endhead
Зрение & Настроить масштаб IDE и браузера, избегать бликов, делать паузы для глаз & Снижение утомления и сухости глаз\\ \hline
Посадка & Отрегулировать кресло, держать спину с опорой, расположить монитор перед собой & Уменьшение статической нагрузки\\ \hline
Кисти и плечи & Размещать клавиатуру и мышь на одной высоте, использовать горячие клавиши, делать разминку & Снижение нагрузки на руки\\ \hline
Рабочее место & Убрать лишние предметы, организовать кабели, не размещать жидкости рядом с техникой & Снижение бытовых и электрических рисков\\ \hline
Режим работы & Делать короткие перерывы каждые 45-60 минут, чередовать код, чтение и проверку интерфейса & Поддержание концентрации\\ \hline
Микроклимат & Проветривать помещение, поддерживать комфортную температуру и влажность & Снижение усталости и сонливости\\ \hline
Планирование & Использовать чек-листы запуска, сборки и проверки отчета & Уменьшение стресса и количества ошибок\\ \hline
\end{longtable}
\normalsize


Для дипломного проекта особенно полезен чек-лист завершения работы. В него входят запуск backend, проверка frontend, сборка LaTeX-отчета, проверка наличия рисунков, просмотр содержания и контроль списка источников. Такой чек-лист уменьшает вероятность случайно сдать неполный комплект материалов.


Рациональная организация труда влияет не только на безопасность, но и на качество результата. Если разработчик работает без перерывов, чаще появляются невнимательные ошибки: пропущенные зависимости, неправильные переменные окружения, несобранный PDF или устаревшие скриншоты. Поэтому охрана труда в данном проекте связана с инженерной дисциплиной и общей надежностью подготовки дипломной работы.


\section{Организация безопасного завершения работы}

Завершение работы за компьютером должно быть таким же организованным, как и начало работы. После длительной разработки важно не только сохранить файлы, но и убедиться, что рабочее место остается безопасным: оборудование не перегревается, зарядные устройства отключены при необходимости, кабели не находятся под нагрузкой, а на столе нет жидкости рядом с техникой.


Для дипломного проекта безопасное завершение работы также включает контроль цифровых материалов. Если отчет собирается из LaTeX, нужно проверить, что финальный PDF создан после последних изменений, а изображения находятся в папке проекта и не потеряны при переносе архива. Если проект демонстрируется на защите, необходимо заранее убедиться, что backend, frontend и база данных запускаются в выбранном сценарии.


\noindent\scriptsize
\setlength{\tabcolsep}{3pt}
\renewcommand{\arraystretch}{1.18}
\begin{longtable}{|>{\RaggedRight\arraybackslash}p{0.24\linewidth}|>{\RaggedRight\arraybackslash}p{0.34\linewidth}|>{\RaggedRight\arraybackslash}p{0.24\linewidth}|}
\caption{Контрольный список завершения работы}\label{tab:3-3}\\
\hline
Этап & Действие & Назначение\\ \hline
\endfirsthead
\hline
Этап & Действие & Назначение\\ \hline
\endhead
Сохранение материалов & Закрыть редакторы после сохранения исходных файлов, отчета и скриншотов & Предотвращение потери результатов работы\\ \hline
Проверка сборки & Собрать финальный PDF и убедиться, что отсутствуют ошибки компиляции & Контроль готовности пояснительной записки\\ \hline
Проверка демонстрации & Запустить основные команды проекта или зафиксировать проверенный сценарий запуска & Снижение риска сбоя на защите\\ \hline
Оборудование & Корректно завершить работу приложений, выключить или перевести устройство в безопасный режим & Снижение нагрузки на оборудование\\ \hline
Рабочее место & Убрать лишние предметы, проверить кабели, убрать жидкости от техники & Профилактика бытовых и электрических рисков\\ \hline
Самочувствие & Сделать перерыв после длительной сессии, размять кисти и шею, дать отдых глазам & Восстановление работоспособности\\ \hline
\end{longtable}
\normalsize


Применение такого контрольного списка особенно полезно перед сдачей дипломного проекта. Он объединяет требования охраны труда и практическую дисциплину разработки: рабочее место остается безопасным, а комплект материалов не расходится с последней версией проекта. В результате студент снижает как физические, так и организационные риски подготовки к защите.


Перед демонстрацией проекта рекомендуется подготовить рабочее место заранее: зарядить ноутбук, проверить доступ к сети, закрыть лишние приложения, открыть необходимые вкладки и убедиться, что проектные материалы находятся в одной папке. Это снижает спешку перед защитой и уменьшает вероятность ошибочных действий, связанных с усталостью или волнением.


Если защита проходит в аудитории, следует учитывать расположение кабелей питания, проектора и рабочих мест других участников. Кабели не должны проходить через проход без фиксации, устройство не должно перегреваться под плотной крышкой или рядом с источниками тепла, а экран должен быть расположен так, чтобы докладчик мог работать без неудобного наклона корпуса. Эти меры простые, но они поддерживают безопасность и помогают провести демонстрацию без лишних технических задержек.
